\section{Read the results}

\textbf{Today.} The model writes tables of production, consumption, flows, infrastructure investment and prices for every node and period.
The analysts check them for consistency and plausibility, and an implausible result can send them back to an earlier step: in the working session, the example was a short export boom that the model could not justify building the infrastructure for, which led to a correction of the storyline rather than of the model.
The results are then interpreted as findings, in journal articles such as the analysis of whether the European Union still needs Russian gas \citep{barner2025russian} and in policy-facing reports \citep{holz2024diw}.
This step was discussed least in the working session, so its description rests mainly on the documentation and the published studies.

\textbf{What the step needs.} Three things, again of different kinds.
Comparing scenarios, computing balances and totals, and tracing which flows and investments changed is arithmetic.
Explaining why a result looks as it does, for instance why prices in one region rise more than in another when a route closes, is reasoning across the model's structure and the scenario's inputs.
Deciding what the results mean, and what may be claimed from them, is the research itself.

\textbf{What an agent could do.} Code should produce the comparisons: differences between scenarios, utilisation of each arc, and the deviations from the reference values that Section 5.5 already requires.
An agent could then draft explanations, linking a change in the results to the inputs that changed, and flag patterns that experienced analysts would question, such as capacity that is built and then barely used, or a trade flow that reverses from one period to the next.
Explaining an energy model's results to people who did not build it is one of the uses of LLMs already proposed for energy-system modelling \citep{forootani2025re}.
As in the systems of Section 2.4, every number in such an explanation must come from the result tables, never from the model's own text.

\textbf{Where a person stays involved.} The interpretation stays with the analysts.
Two weaknesses from Section 4.3 make this more than a formality.
An LLM tends to agree with the view it is presented with, so an agent that knows the storyline may explain away a result that should have prompted a revision.
And a fluent explanation raises a reader's confidence whether or not it is right, so the agent's drafts are best treated as a list of hypotheses for the analyst to check, each pointing to the tables that support it.
